\section{Introduction}
A persistent-memory agent can change internally long before task accuracy reveals anything. Consider two copies of the same interaction trajectory. If one is labeled successful and the other failed before reflection, an outcome-conditioned writer can store different reusable advice although the underlying experience is fixed. That establishes a persistent-state intervention, but not behavioral self-improvement: the different state must still be exposed by retrieval, influence a decision, and survive to an outcome.

We study this gap as \emph{stage-resolved transport}. For a frozen trajectory $\tau$, we switch the released success/failure reflection branch---a bundled writer-protocol intervention that changes reflection instruction together with outcome/reward semantics---and instrument
\[
\text{write }(W)\rightarrow\text{source-item exposure }(E)\rightarrow\text{first-action uptake }(U)\rightarrow\text{outcome }(O).
\]
A separate forced fixed-evidence arm $F$ bypasses native retrieval and asks whether the paired memory content can affect the downstream system when supplied. These surfaces answer different questions and have different units. A write difference proves state change, a retrieval hit proves availability of the source item, and neither proves that the branch-differentiating residual changed the policy. Conversely, a weak endpoint alone cannot tell whether the update failed to write, failed to retrieve, was ignored after retrieval, or changed behavior only sparsely. We call this loss of mechanism information \emph{evaluation coarsening}.

The frozen evidence produces exactly this disagreement across surfaces. Reward-conditioned writing diverges for 20/20 Shopping pairs and 4/4 Reddit pairs; a stronger same-mode wording control leaves a reward-mode excess of 0.105 ($p=0.0078$). Forced supply yields terminal $|\Delta|=0.15625$ ($p=0.00074$). Yet native Shopping retrieves the branch-conditioned source item in 125/172 held-out opportunities while the first-action branch contrast is not supported (TV $=0.06944$, $p=0.5801$, 0/36 modal changes) and terminal transport is sparse ($|\Delta|=0.02083$, $p=0.4289$, 34/36 zero). Reddit again has reliable write divergence but 6/8 zero terminal cells and opposite signs in the two nonzero cells. The strongest supported operational localization is therefore after source-item exposure and before stable first-action uptake. It is an evidence boundary, not a causal attenuation onset or mediation coefficient.

Stage resolution is not merely ``report more metrics.'' The native prerequisite order $O\Rightarrow U\Rightarrow E\Rightarrow W$ permits five binary prefix depths; crossing them with $F\in\{0,1\}$ yields a complete 10-state diagnostic class for this experiment. Exhaustive enumeration of all 32 observed-surface subsets shows that $W/E/U/O/F$ is the unique separating basis: omitting a native coordinate aliases adjacent prefix states, while omitting $F$ aliases the two capacity states at every depth. This paper-specific result explains why coarsening can make genuinely different failure explanations observationally identical.

The diagnostic also creates an engineering question: does the first unsupported measured stage identify a useful repair family? We test this prospectively in two steps rather than assuming a repair. First, on 13 outcome-blind states we compare native prompting with matched decision-check controls; an explicit reminder to use memory leaves mean uptake unchanged ($U_{A0}=U_{A1}=U_{A2}=0.0962$; 312/312 calls complete), so salience alone does not repair the boundary. Second, on a disjoint fresh 12-state pilot we compare native memory, a memory-only adapter, and a state-conditioned binding adapter that converts the same retrieved memory into a concise current-state action implication before policy execution. The adapter is deliberately a proof-of-concept, not a novel memory architecture. Mean uptake rises to $0.208$ from $0.125$ natively and $0.083$ under memory-only adaptation, with mean $D=U_{A2}-U_{A1}=0.125$. The effect is nevertheless concentrated: $D>0$ in only 3/12 states, 4 are negative, and the preregistered cross-state consistency gate fails. The combined result is sharper than either pilot alone: \emph{the diagnosis selects a mechanism-aligned repair family and state binding can reopen transport in some states, but repair authority remains conditional rather than automatic}.

Our contributions are:
\begin{enumerate}
    \item \textbf{Controlled stage-resolved transport.} A same-trajectory success/failure writer intervention follows one induced persistent-state contrast through native exposure, first-action uptake, outcome, and an explicit forced-capacity bypass control.
    \item \textbf{Diagnostic completeness under the experiment's prerequisite class.} The complete 10-state class induced by $O\Rightarrow U\Rightarrow E\Rightarrow W$ and $F$ requires the full $W/E/U/O/F$ observation basis to avoid aliases.
    \item \textbf{An empirical transport boundary.} Write divergence and source-item exposure are established while stable first-action uptake is not, and cross-domain terminal effects are sparse and sign-heterogeneous.
    \item \textbf{Two prospective repair probes that separate salience from state-conditioned binding.} A simple explicit-use reminder has no aggregate effect, while a fresh state-conditioned adapter increases mean first-action uptake relative to native and memory-only controls but fails the preregistered cross-state consistency gate. This establishes a heterogeneous proof-of-concept opening without promoting the adapter to a universal repair or a new method claim.
\end{enumerate}

The resulting claim is deliberately narrower than a new memory-utilization method: \textbf{state divergence is not behavioral authority; availability is not use; diagnosis is not a repair prescription}.
